% ======================================================================
% DRAFT strengthening content for review + placement. Not yet in main_v1.
% (1) frontier-comparison table  -- suggested: Introduction or Related Work
% (2) LLM positioning paragraph  -- suggested: end of Introduction, or a
%     short Discussion paragraph before Limitations
% ======================================================================

% ---------- (1) frontier-comparison table ----------------------------
\begin{table}[h]
\centering
\small
\begin{tabular}{lll}
\toprule
work & method & equation class \\
\midrule
Poesia et al.\ (2021)      & contrastive RL          & linear / elementary \\
Dabelow \& Ueda (2024)     & Q-learning, symbolic    & linear \\
\textbf{This work}         & PPO + TreeMLP           & \textbf{nonlinear + open (CoV)} \\
\bottomrule
\end{tabular}
\caption{\textbf{The frontier.} Prior reinforcement-learning / trace-free
symbolic equation solving stops at linear equations. This work extends it to
nonlinear equations (radicals, exponentials, trigonometric functions) and to
\emph{open} equations, which require a generative change of variables.}
\label{tab:frontier}
\end{table}

% ---------- (2) LLM positioning paragraph ----------------------------
\textbf{Relation to large language models.}
Large language models can produce answers to equations of the kind we study,
and we do not position this work as a competitor on solution accuracy. The
contribution is orthogonal. We ask whether an agent can \emph{learn} -- with
no supervised solution traces -- the exact symbolic transformations that
solve an equation step by step. Every action in our setting is an exact
algebraic operation applied by a computer-algebra backend, so a solution is a
\emph{verified derivation} rather than a generated string: the agent cannot
hallucinate an intermediate step, and a claimed solution either checks or it
does not. This exact, trace-free learning setting is precisely the one
studied by \citet{poesia2021contrastive} and \citet{dabelow2024symbolic};
our contribution is to carry it from linear equations to nonlinear and
generative equation solving. Whether a learned symbolic policy of this kind
can be distilled into, or used to guide, a general language model is an
interesting question we leave to future work.
